%==============================
\section{Probabilistic Temporal Tasks and Safe-Return Constraints}
\label{ch06:sec:ch06_safe_plan}

This section considers uncertain motion planning for a robot modeled by a
labeled Markov decision process, abbreviated as MDP. The nominal mission is
specified by an LTL formula, while safety is represented by a separate
safe-return requirement. This separation is needed because productive operation
and permanent residence in a safe set are generally incompatible on the same
infinite execution. The temporal-logic background follows
Chapter~\ref{chap:preliminaries}, and remains
consistent with probabilistic temporal-planning models
\citep{guo2018probabilistic,belta2017formal,ding2014ltl}.

%==============================================================
\subsection{Uncertain system and probabilistic temporal task}
\label{ch06:sec:problem-formulate}

This subsection defines the labeled MDP, the probabilistic temporal-task
specification, and the safe-return requirement that couples task fulfillment
with recoverability.

\subsubsection{Uncertain system as labeled MDP}
\label{ch06:subsec:probMDP}

The robot and its environment are modeled by a finite labeled MDP. This model
captures the stochastic state evolution together with the propositions needed
for temporal-task evaluation. The formal definition is given by:
\begin{equation}
\label{ch06:eq:mdp}
\mathcal{M} \triangleq
\bigl(X, U, D, p_D, AP, L, c_D, x_0\bigr),
\end{equation}
where $X$ is a finite state space, $U(x)$ is the set of admissible control
actions at state $x \in X$, and
$D \triangleq \{(x,u) \mid x \in X,\; u \in U(x)\}$ is the admissible
state-action set. The transition probability function
$p_D \colon X \times U \times X \rightarrow [0,1]$ satisfies
$\sum_{\check{x} \in X} p_D(x,u,\check{x}) = 1$ for every $(x,u) \in D$.
Moreover, $AP$ is the set of atomic propositions, the labeling map
$L \colon X \rightarrow 2^{AP}$ assigns the propositions satisfied at each
state, $c_D \colon D \rightarrow \mathbb{R}_{>0}$ is the stage-cost function,
and $x_0 \in X$ is the initial state. Equation~\eqref{ch06:eq:mdp} follows the
standard MDP model of \citet{puterman2014markov}, augmented with a deterministic
labeling map.
The deterministic labeling in \eqref{ch06:eq:mdp} is sufficient below. If
workspace features are uncertain, the construction can be extended by state
augmentation or by probabilistic labeling
\citep{guo2018probabilistic}. For instance, in search and rescue, states may
encode regions and local conditions, actions may encode motion primitives, and
labels may mark exits, victims, charging stations, or hazards.


\subsubsection{Task specification}
\label{ch06:subsec:LTL}


The nominal mission is an LTL formula $\varphi$ over the proposition set
$AP$ from \eqref{ch06:eq:mdp}. The syntax and semantics follow
Section~\ref{ch02:sec:LTL}; the automata-based construction follows
Sections~\ref{ch02:subsec:NBA}--\ref{ch02:subsec:product}, using
deterministic automata for probabilistic planning
\citep{baier2008principles,belta2017formal}.
Typical tasks include persistent surveillance,
$\varphi = (\Box \Diamond r_1) \wedge (\Box \Diamond r_2)
\wedge (\Box \Diamond r_3)$, ordered service,
$\varphi = \Diamond \bigl((\texttt{pick} \wedge r_1) \wedge
\Diamond (\texttt{drop} \wedge r_2)\bigr)$, and response requirements,
$\varphi = \Box \bigl(\texttt{low} \rightarrow
\Diamond (\texttt{supply} \wedge r_1)\bigr)$.


For each stage $T \geq 0$, let
$R_T \triangleq x_0 l_0 u_0 \cdots x_T l_T$ denote a finite run, where
$l_t = L(x_t)$ and $p_D(x_t,u_t,x_{t+1}) > 0$ for every valid transition.
Let $\mathbf{R}_T$ be the set of all runs of length $T+1$, and let
$\mathbf{R}_{\infty}$ be the set of all infinite runs. A finite-memory policy
is written as $\boldsymbol{\mu} \triangleq \mu_0 \mu_1 \cdots$, where each
$\mu_t \colon \mathbf{R}_t \times U \rightarrow [0,1]$ assigns a distribution
over admissible actions based on the observed history. The set of such policies
is denoted by $\overline{\boldsymbol{\mu}}$.
Given $\boldsymbol{\mu} \in \overline{\boldsymbol{\mu}}$, the MDP in
\eqref{ch06:eq:mdp} induces a probability measure over
$\mathbf{R}_{\infty}$ \citep{baier2008principles}. The corresponding task
satisfaction probability is given by:
\begin{equation}
\label{ch06:eq:satisfy}
\textbf{Sat}_{\mathcal{M}}^{\boldsymbol{\mu}} \triangleq
Pr_{\mathcal{M}}^{\boldsymbol{\mu}}
\bigl(
\{R_{\infty} \in \mathbf{R}_{\mathcal{M}}^{\boldsymbol{\mu}} \mid
R_{\infty}|_{L} \models \varphi\}
\bigr),
\end{equation}
where $\mathbf{R}_{\mathcal{M}}^{\boldsymbol{\mu}} \subseteq
\mathbf{R}_{\infty}$ is the set of infinite runs generated by
$\mathcal{M}$ under $\boldsymbol{\mu}$. Equation~\eqref{ch06:eq:satisfy} is
therefore the probability that the induced label sequence satisfies the task
formula.
The outbound performance criterion is the expected mean cost:
\begin{equation}
\label{ch06:eq:total-cost}
\textbf{Cost}_{\mathcal{M}}^{\boldsymbol{\mu}} \triangleq
\mathbb{E}_{\mathcal{M}}^{\boldsymbol{\mu}}
\left[
\limsup_{T \rightarrow \infty}
\frac{1}{T} \sum_{t=0}^{T-1} c_D(x_t,u_t)
\right],
\end{equation}
which is suited to persistent tasks over infinite executions. An average-cost
criterion is used instead of a stochastic-shortest-path cost, since the latter
is tailored to terminating tasks
\citep{puterman2014markov,bertsekas1995neuro,polychronopoulos1996stochastic}.

\subsubsection{Safe-return constraints}
\label{ch06:subsec:safety}

Safe return is described by a separate LTL formula requiring that, after a
return request, the robot reaches a designated safe set and remains there:
\begin{equation}
\label{ch06:eq:safe-task}
\varphi_{\texttt{r}} \triangleq
\varphi_{\texttt{r}}^{1} \wedge \Diamond \Box \,\varphi_{\texttt{r}}^{2},
\end{equation}
where $\varphi_{\texttt{r}}^{1}$ is an sc-LTL formula describing the transient
return phase, and $\varphi_{\texttt{r}}^{2}$ identifies the terminal safe set.
More precisely,
$\varphi_{\texttt{r}}^{2} \triangleq \bigvee_{l \in L_s} l$, where
$L_s \subseteq AP$ is the set of labels associated with the safe states.
Equation~\eqref{ch06:eq:safe-task} requires eventual completion of the transient
phase and, afterwards, permanent residence in the safe set
\citep{baier2008principles,pnueli1981temporal}.

Under a return policy
$\boldsymbol{\mu}_{\texttt{r}} \in \overline{\boldsymbol{\mu}}$, the
probability of satisfying the safe-return specification in
\eqref{ch06:eq:safe-task} is given by:
\begin{equation}
\label{ch06:eq:sat-return}
\textbf{Sat}_{\mathcal{M}}^{\boldsymbol{\mu}_{\texttt{r}}} \triangleq
Pr_{\mathcal{M}}^{\boldsymbol{\mu}_{\texttt{r}}}
\bigl(
\{R_{\infty} \in \mathbf{R}_{\mathcal{M}}^{\boldsymbol{\mu}_{\texttt{r}}}
\mid R_{\infty}|_{L} \models \varphi_{\texttt{r}}\}
\bigr),
\end{equation}
where $\mathbf{R}_{\mathcal{M}}^{\boldsymbol{\mu}_{\texttt{r}}}$ denotes the
set of infinite runs generated under the return policy.
For a search-and-rescue mission, a representative safe-return requirement is
$\varphi_{\texttt{r}} = \Diamond \texttt{ex} \wedge
\Diamond \Box \,\texttt{bs}$, which requires the robot to pass through a
designated exit and then remain at the base station.
The formula $\varphi_{\texttt{r}}$ differs from the nominal task $\varphi$.
Collision avoidance and emergency-response requirements are encoded in
$\varphi$, whereas $\varphi_{\texttt{r}}$ expresses recoverability: at any
externally chosen switching time, execution should admit a return policy with
sufficient success probability. This interpretation is related to safe
exploration and recoverability constraints
\citep{moldovan2012safe,turchetta2016safe}, but is coupled here with a general
temporal mission.
Let $\boldsymbol{\mu}_{\texttt{o}}$ denote the outbound policy for the nominal
task, and let $\boldsymbol{\mu}_{\texttt{r}}$ denote the return policy
activated after a return request. Let $\mathcal{M}'_{\texttt{o}}$ denote the
same MDP as in \eqref{ch06:eq:mdp}, but initialized with the state
distribution induced at the switching instant under
$\boldsymbol{\mu}_{\texttt{o}}$. The safety level of the pair
$(\boldsymbol{\mu}_{\texttt{o}}, \boldsymbol{\mu}_{\texttt{r}})$ is then
defined by:
\begin{equation}
\label{ch06:eq:sat-return-new}
\textbf{Safe}_{\mathcal{M}}^{
\boldsymbol{\mu}_{\texttt{o}},\boldsymbol{\mu}_{\texttt{r}}
}
\triangleq
\textbf{Sat}_{\mathcal{M}'_{\texttt{o}}}^{
\boldsymbol{\mu}_{\texttt{r}}
},
\end{equation}
so \eqref{ch06:eq:sat-return-new} measures the return success probability from
the distribution induced by the outbound execution.

\begin{definition}
\label{ch06:def:safety}
Given the MDP $\mathcal{M}$ and an outbound policy
$\boldsymbol{\mu}_{\texttt{o}}$, the policy
$\boldsymbol{\mu}_{\texttt{o}}$ is called $\chi_{\texttt{r}}$-safe if there
exists a return policy $\boldsymbol{\mu}_{\texttt{r}}$ such that:
\begin{equation}
\label{ch06:eq:safe}
\textbf{Safe}_{\mathcal{M}}^{
\boldsymbol{\mu}_{\texttt{o}},\boldsymbol{\mu}_{\texttt{r}}
}
\geq \chi_{\texttt{r}},
\end{equation}
where $\chi_{\texttt{r}} \in [0,1]$ is the prescribed safety threshold.
\end{definition}

\subsubsection{Problem statement}
\label{ch06:subsec:problem}

\begin{problem}
\label{ch06:prob:main}
Given the labeled MDP $\mathcal{M}$ in \eqref{ch06:eq:mdp}, the task formula
$\varphi$, and the safe-return formula $\varphi_{\texttt{r}}$, synthesize an
outbound policy $\boldsymbol{\mu}_{\texttt{o}}$ and a return policy
$\boldsymbol{\mu}_{\texttt{r}}$ such that
\begin{equation}
\label{ch06:eq:objective}
\begin{aligned}
\min_{
\boldsymbol{\mu}_{\texttt{o}},\boldsymbol{\mu}_{\texttt{r}}
\in \overline{\boldsymbol{\mu}}
}
\quad
&\textbf{Cost}_{\mathcal{M}}^{\boldsymbol{\mu}_{\texttt{o}}}, \\
\text{s.t.}\quad
&\textbf{Sat}_{\mathcal{M}}^{\boldsymbol{\mu}_{\texttt{o}}}
\geq \chi_{\texttt{o}}, \; \textbf{Safe}_{\mathcal{M}}^{
\boldsymbol{\mu}_{\texttt{o}},\boldsymbol{\mu}_{\texttt{r}}
}
\geq \chi_{\texttt{r}};
\end{aligned}
\end{equation}
where $\chi_{\texttt{o}}, \chi_{\texttt{r}} \in [0,1]$ are prescribed lower
bounds on task satisfaction and safe return.
\end{problem}

Problem~\ref{ch06:prob:main} above is a two-policy synthesis problem rather than a
single-policy multi-objective problem. The outbound policy is admissible only
if it preserves a return policy satisfying the safety constraint in~\eqref{ch06:eq:safe}.


%=========================================================
%=========================================================
%=========================================================
\subsection{Product-based reformulation and sequential synthesis}
\label{ch06:sec:analysis}

This subsection reformulates the two constraints in
Problem~\ref{ch06:prob:main} on product automata. The return policy is computed
first, and the outbound policy is then synthesized under the induced
safe-return values. The construction follows the standard automata-based
treatment of MDPs under temporal logic specifications
\citep{klein2007ltl2dstar,baier2008principles,belta2017formal,
guo2018probabilistic}.

\subsubsection{Product automata and accepting components}
\label{ch06:subsec:product}

To connect temporal specifications with the MDP model, deterministic Rabin
automata are used. Given an LTL formula \(\psi\) over the proposition set
\(AP\), let \(\mathcal{A}_{\psi}\) denote the associated deterministic Rabin
automaton, written as:
\begin{equation}
\label{ch06:eq:dra}
\mathcal{A}_{\psi} \triangleq
\bigl(
Q_{\psi}, 2^{AP}, \delta_{\psi}, q_{\psi}^{0},
\mathrm{Acc}_{\mathcal{A}_{\psi}}
\bigr),
\end{equation}
where \(Q_{\psi}\) is a finite state set, \(2^{AP}\) is the input alphabet,
\(\delta_{\psi} \colon Q_{\psi} \times 2^{AP} \rightarrow Q_{\psi}\) is the
transition function, \(q_{\psi}^{0} \in Q_{\psi}\) is the initial state, and
\(\mathrm{Acc}_{\mathcal{A}_{\psi}}\) is the Rabin accepting condition.
The accepting condition is given by:
\begin{equation}
\label{ch06:eq:dra-acc}
\mathrm{Acc}_{\mathcal{A}_{\psi}} \triangleq
\left\{
\bigl(H_{\mathcal{A}_{\psi}}^{1}, I_{\mathcal{A}_{\psi}}^{1}\bigr),
\cdots,
\bigl(H_{\mathcal{A}_{\psi}}^{N}, I_{\mathcal{A}_{\psi}}^{N}\bigr)
\right\},
\end{equation}
where
\(H_{\mathcal{A}_{\psi}}^{i}, I_{\mathcal{A}_{\psi}}^{i} \subseteq Q_{\psi}\)
for each \(i = 1,\cdots,N\). Given an infinite word
\(w = l_{0} l_{1} l_{2} \cdots\) over \(2^{AP}\), the corresponding run of
\(\mathcal{A}_{\psi}\) is the sequence
\(q_{0} q_{1} q_{2} \cdots\), where \(q_{0} = q_{\psi}^{0}\) and
\(q_{k+1} = \delta_{\psi}(q_{k}, l_{k})\) for all \(k \geq 0\). This run is
accepting if there exists at least one pair
\(\bigl(H_{\mathcal{A}_{\psi}}^{i}, I_{\mathcal{A}_{\psi}}^{i}\bigr)\) such
that \(H_{\mathcal{A}_{\psi}}^{i}\) is visited only finitely many times and
\(I_{\mathcal{A}_{\psi}}^{i}\) is visited infinitely many times
\citep{baier2008principles}. Translation tools such as \texttt{ltl2dstar} can
be used to construct \(\mathcal{A}_{\psi}\) from \(\psi\)
\citep{klein2007ltl2dstar}.
For the task formula \(\varphi\), let
\begin{equation}
\label{ch06:eq:dra-task}
\mathcal{A}_{\textup{\texttt{o}}} \triangleq
\bigl(
Q_{\textup{\texttt{o}}},
2^{AP},
\delta_{\textup{\texttt{o}}},
q_{\textup{\texttt{o}}}^{0},
\mathrm{Acc}_{\mathcal{A}_{\textup{\texttt{o}}}}
\bigr),
\end{equation}
be the associated deterministic Rabin automaton. The product automaton between
\(\mathcal{M}\) in \eqref{ch06:eq:mdp} and
\(\mathcal{A}_{\textup{\texttt{o}}}\) is defined as follows.

\begin{definition}
\label{ch06:def:product-o}
The product automaton for the nominal task is
\(\mathcal{P}_{\textup{\texttt{o}}} \triangleq
\mathcal{M} \times \mathcal{A}_{\textup{\texttt{o}}}\), given by:
\begin{equation}
\label{ch06:eq:prod-o}
\mathcal{P}_{\textup{\texttt{o}}} \triangleq
\bigl(
S_{\textup{\texttt{o}}},
U_{\textup{\texttt{o}}},
E_{\textup{\texttt{o}}},
p_{E,\textup{\texttt{o}}},
c_{E,\textup{\texttt{o}}},
s^{0}_{\textup{\texttt{o}}},
\mathrm{Acc}_{\mathcal{P}_{\textup{\texttt{o}}}}
\bigr),
\end{equation}
where \(S_{\textup{\texttt{o}}} \subseteq
X \times Q_{\textup{\texttt{o}}}\) is the reachable product-state set,
\(U_{\textup{\texttt{o}}}(\langle x,q\rangle) \triangleq U(x)\), and
\(E_{\textup{\texttt{o}}} \triangleq
\{(s,u) \mid s \in S_{\textup{\texttt{o}}},
u \in U_{\textup{\texttt{o}}}(s)\}\).
The transition probability is given by:
\begin{equation}
\label{ch06:eq:new_pro}
p_{E,\textup{\texttt{o}}}
\bigl(
\langle x,q\rangle, u, \langle \check{x},\check{q}\rangle
\bigr)
\triangleq p_D(x,u,\check{x}),
\end{equation}
where \((x,u) \in D\) and
\(\check{q} = \delta_{\textup{\texttt{o}}}(q,L(x))\).
The stage cost is
\(c_{E,\textup{\texttt{o}}}(\langle x,q\rangle,u)
\triangleq c_D(x,u)\), the initial state
\(s^{0}_{\textup{\texttt{o}}} \triangleq
\langle x_0, q^{0}_{\textup{\texttt{o}}}\rangle\), and the accepting pairs
induced from
\(\mathrm{Acc}_{\mathcal{A}_{\textup{\texttt{o}}}}\) in the standard
manner \citep{baier2008principles}.
\end{definition}

The product automaton \(\mathcal{P}_{\textup{\texttt{o}}}\) combines traces
of \(\mathcal{M}\) with words accepted by
\(\mathcal{A}_{\textup{\texttt{o}}}\). Its accepting behavior is captured by
accepting maximum end components, abbreviated as AMECs. Let
\(\Xi_{\mathrm{acc}}^{\textup{\texttt{o}}}\) denote the set of AMECs of
\(\mathcal{P}_{\textup{\texttt{o}}}\). If
\((S^{\textup{\texttt{o}}}_{c}, U^{\textup{\texttt{o}}}_{c}) \in
\Xi_{\mathrm{acc}}^{\textup{\texttt{o}}}\), then
\(S^{\textup{\texttt{o}}}_{c}\) is an accepting recurrent region and
\(U^{\textup{\texttt{o}}}_{c}\) specifies the admissible actions that keep
the product inside that region. The union of all accepting components is
given by:
\begin{equation}
\label{ch06:eq:AMECs}
S^{\textup{\texttt{o}}}_{\Xi} \triangleq
\bigcup_{
(S^{\textup{\texttt{o}}}_{c},U^{\textup{\texttt{o}}}_{c})
\in \Xi_{\mathrm{acc}}^{\textup{\texttt{o}}}
}
S^{\textup{\texttt{o}}}_{c},
\end{equation}
which is the set of product states from which the accepting condition can be
sustained indefinitely,
see Definitions~10.116, 10.117, 10.124 of~\citet{baier2008principles}.
Let \(\mathcal{A}_{\textup{\texttt{r}}}\) be the deterministic Rabin
automaton associated with the safe-return formula
\(\varphi_{\texttt{r}}\). The corresponding product automaton is defined in
the same way.
The product automaton for the safe-return specification is
\(\mathcal{P}_{\textup{\texttt{r}}} \triangleq
\mathcal{M} \times \mathcal{A}_{\textup{\texttt{r}}}\). Its construction is
the same as in Definition~\ref{ch06:def:product-o}, with
\(\mathcal{A}_{\textup{\texttt{o}}}\) replaced by
\(\mathcal{A}_{\textup{\texttt{r}}}\). Let
\(\Xi_{\mathrm{acc}}^{\textup{\texttt{r}}}\) denote the set of AMECs of
\(\mathcal{P}_{\textup{\texttt{r}}}\), and define
$S^{\textup{\texttt{r}}}_{\Xi} \triangleq
\bigcup_{
(S^{\textup{\texttt{r}}}_{c},U^{\textup{\texttt{r}}}_{c})
\in \Xi_{\mathrm{acc}}^{\textup{\texttt{r}}}
}S^{\textup{\texttt{r}}}_{c}$.
Since both automata are deterministic, product policies and MDP policies
correspond uniquely. Product policies are denoted by
\(\boldsymbol{\pi}_{\textup{\texttt{o}}}\) and
\(\boldsymbol{\pi}_{\textup{\texttt{r}}}\); the corresponding MDP policies
remain \(\boldsymbol{\mu}_{\textup{\texttt{o}}}\) and
\(\boldsymbol{\mu}_{\textup{\texttt{r}}}\).

\subsubsection{Reformulation of the task and safe-return constraints}
\label{ch06:subsec:task-reformulate}

Task satisfaction becomes a reachability property on the outbound product.
The safe-return constraint is handled similarly, but its value is computed on
the return product and lifted to the outbound product through the switching
state. Theorems~\ref{ch06:theorem:task-bound} and
\ref{ch06:theorem:safe-bound} convert the two constraints in
Problem~\ref{ch06:prob:main} into product-automaton quantities.

\begin{theorem}
\label{ch06:theorem:task-bound}
For any outbound policy
\(\boldsymbol{\mu}_{\textup{\texttt{o}}}\) and its corresponding product
policy \(\boldsymbol{\pi}_{\textup{\texttt{o}}}\), the task-satisfaction
probability in \eqref{ch06:eq:satisfy} satisfies:
\begin{equation}
\label{ch06:eq:theo-sat-bound}
\textbf{Sat}_{\mathcal{M}}^{
\boldsymbol{\mu}_{\textup{\texttt{o}}}}
=
\mathbb{E}_{\widetilde{\mathcal{P}}_{\textup{\texttt{o}}}}^{
\boldsymbol{\pi}_{\textup{\texttt{o}}}}
\left[
\sum_{t=0}^{\infty}
b_{s_t,S^{\textup{\texttt{o}}}_{\Xi}}
\right],
\end{equation}
where \(b_{s,S} \triangleq \mathds{1}_{\{s \in S\}}\), and
\(\widetilde{\mathcal{P}}_{\textup{\texttt{o}}}\) is obtained from
\(\mathcal{P}_{\textup{\texttt{o}}}\) by making
\(S^{\textup{\texttt{o}}}_{\Xi}\) absorbing.
\end{theorem}

\begin{proof}
For deterministic Rabin acceptance, a run satisfies \(\varphi\) if and only
if the corresponding run on \(\mathcal{P}_{\textup{\texttt{o}}}\) reaches an
AMEC and then remains inside an accepting component under an admissible
policy \citep{baier2008principles,forejt2011quantitative,
guo2018probabilistic}. Hence, the probability of satisfying \(\varphi\)
equals the probability of reaching
\(S^{\textup{\texttt{o}}}_{\Xi}\) at least once. In the modified product
\(\widetilde{\mathcal{P}}_{\textup{\texttt{o}}}\), the set
\(S^{\textup{\texttt{o}}}_{\Xi}\) is absorbing, so the indicator
\(b_{s_t,S^{\textup{\texttt{o}}}_{\Xi}}\) is counted exactly once along each
accepting trajectory. This yields \eqref{ch06:eq:theo-sat-bound}.
\end{proof}

\begin{theorem}
\label{ch06:theorem:safe-bound}
For any pair of policies
\((\boldsymbol{\mu}_{\textup{\texttt{o}}},
\boldsymbol{\mu}_{\textup{\texttt{r}}})\) and the corresponding product
policies
\((\boldsymbol{\pi}_{\textup{\texttt{o}}},
\boldsymbol{\pi}_{\textup{\texttt{r}}})\), the safe-return quantity in
\eqref{ch06:eq:safe} can be written as
\begin{equation}
\label{ch06:eq:theo-safe-bound}
\textbf{Safe}_{\mathcal{M}}^{
\boldsymbol{\mu}_{\textup{\texttt{o}}},
\boldsymbol{\mu}_{\textup{\texttt{r}}}}
=
\mathbb{E}_{\mathcal{P}_{\textup{\texttt{o}}}}^{
\boldsymbol{\pi}_{\textup{\texttt{o}}}}
\left[
\sum_{t=0}^{\infty}
v^{\star}_{\textup{\texttt{r}}}
\bigl(s^{\textup{\texttt{r}}}_{t}\bigr)
\right],
\end{equation}
where
\(s^{\textup{\texttt{r}}}_{t} \triangleq
\langle x_t, q^{0}_{\textup{\texttt{r}}}\rangle\) is the product state in
\(\mathcal{P}_{\textup{\texttt{r}}}\) associated with the MDP state \(x_t\),
and
\begin{equation}
\label{ch06:eq:vrstar}
v^{\star}_{\textup{\texttt{r}}}(s) \triangleq
\mathbb{E}_{\widetilde{\mathcal{P}}_{\textup{\texttt{r}}}}^{
s,\boldsymbol{\pi}_{\textup{\texttt{r}}}}
\left[
\sum_{\ell=0}^{\infty}
b_{s_{\ell},S^{\textup{\texttt{r}}}_{\Xi}}
\right],
\end{equation}
where \(\widetilde{\mathcal{P}}_{\textup{\texttt{r}}}\) is obtained from
\(\mathcal{P}_{\textup{\texttt{r}}}\) by making
\(S^{\textup{\texttt{r}}}_{\Xi}\) absorbing.
\end{theorem}

\begin{proof}
Consider an arbitrary switching stage \(t \geq 0\). Conditional on the
state \(x_t\) reached under the outbound policy at stage \(t\), the
subsequent evolution of the system is governed by the return policy.
For the modified initial state \(x_t\), the probability of
satisfying \(\varphi_{\texttt{r}}\) under
\(\boldsymbol{\mu}_{\textup{\texttt{r}}}\) is equal to the probability of
reaching \(S^{\textup{\texttt{r}}}_{\Xi}\) in the return product. By the
same argument used in Theorem~\ref{ch06:theorem:task-bound}, this quantity is
given by \(v^{\star}_{\textup{\texttt{r}}}
(s^{\textup{\texttt{r}}}_{t})\) in \eqref{ch06:eq:vrstar}. Averaging these
values over the outbound evolution yields
\eqref{ch06:eq:theo-safe-bound}.
\end{proof}


\subsubsection{Sequential baseline method}
\label{ch06:subsec:complexity}

The reformulation above yields a sequential baseline method. The return
policy is synthesized first, because its value function enters the outbound
optimization through \eqref{ch06:eq:theo-safe-bound}. The return problem is
formulated as:
\begin{equation}
\label{ch06:eq:new-problem-safe}
\max_{
\boldsymbol{\pi}_{\textup{\texttt{r}}}
\in \overline{\boldsymbol{\pi}}
}
\mathbb{E}_{\widetilde{\mathcal{P}}_{\textup{\texttt{r}}}}^{
\boldsymbol{\pi}_{\textup{\texttt{r}}}}
\left[
\sum_{\ell=0}^{\infty}
b_{s_{\ell},S^{\textup{\texttt{r}}}_{\Xi}}
\right],
\end{equation}
which maximizes the reachability of
\(S^{\textup{\texttt{r}}}_{\Xi}\) in
\(\widetilde{\mathcal{P}}_{\textup{\texttt{r}}}\). Once
\(\boldsymbol{\pi}_{\textup{\texttt{r}}}\) is fixed, the value function
\(v^{\star}_{\textup{\texttt{r}}}\) in \eqref{ch06:eq:vrstar} is computed on
\(\widetilde{\mathcal{P}}_{\textup{\texttt{r}}}\).
The outbound policy is then obtained from the constrained problem below:
\begin{equation}
\label{ch06:eq:new-problem}
\begin{aligned}
\min_{
\boldsymbol{\pi}_{\textup{\texttt{o}}}
\in \overline{\boldsymbol{\pi}}
}
\quad
&\textbf{Cost}_{\mathcal{P}_{\textup{\texttt{o}}}}^{
\boldsymbol{\pi}_{\textup{\texttt{o}}}}, \\
\text{s.t.}\quad
&\mathbb{E}_{\widetilde{\mathcal{P}}_{\textup{\texttt{o}}}}^{
\boldsymbol{\pi}_{\textup{\texttt{o}}}}
\left[
\sum_{t=0}^{\infty}
b_{s_t,S^{\textup{\texttt{o}}}_{\Xi}}
\right]
\geq \chi_{\texttt{o}}, \quad
\mathbb{E}_{\mathcal{P}_{\textup{\texttt{o}}}}^{
\boldsymbol{\pi}_{\textup{\texttt{o}}}}
\left[
\sum_{t=0}^{\infty}
v^{\star}_{\textup{\texttt{r}}}
\bigl(s^{\textup{\texttt{r}}}_{t}\bigr)
\right]
\geq \chi_{\texttt{r}};
\end{aligned}
\end{equation}
where the constraints follow from \eqref{ch06:eq:theo-sat-bound} and
\eqref{ch06:eq:theo-safe-bound}. Equation~\eqref{ch06:eq:new-problem} is a
constrained MDP problem
\citep{altman1999constrained,bertsekas1991analysis}; its prefix and suffix
policies can be synthesized by an occupancy-measure linear program, as in
probabilistic temporal-logic planning
\citep{guo2018probabilistic,ding2014optimal}. The procedure is summarized in
Algorithm~\ref{ch06:alg:direct}. Although the linear programs are polynomial
in the product sizes, the products grow rapidly with the MDP and automata
sizes. Thus, the direct method is used as a baseline for the structured
synthesis below.

\begin{algorithm}[t]
\caption{Sequential baseline method}
\label{ch06:alg:direct}
\LinesNumbered
\KwIn{\(\mathcal{M}\), \((\varphi,\chi_{\texttt{o}})\), and
\((\varphi_{\texttt{r}},\chi_{\texttt{r}})\).}
\KwOut{\((\mathcal{P}_{\textup{\texttt{r}}},\boldsymbol{\pi}_{\textup{\texttt{r}}},
v^{\star}_{\textup{\texttt{r}}})\), and
\((\mathcal{P}_{\textup{\texttt{o}}},\boldsymbol{\pi}_{\textup{\texttt{o}}})\).}
Construct \(\mathcal{P}_{\textup{\texttt{r}}}\) and compute
\(S^{\textup{\texttt{r}}}_{\Xi}\)\;
Solve \eqref{ch06:eq:new-problem-safe} for
\(\boldsymbol{\pi}_{\textup{\texttt{r}}}\) and
\(v^{\star}_{\textup{\texttt{r}}}\)\;
Construct \(\mathcal{P}_{\textup{\texttt{o}}}\) and compute
\(S^{\textup{\texttt{o}}}_{\Xi}\)\;
Solve \eqref{ch06:eq:new-problem} for
\(\boldsymbol{\pi}_{\textup{\texttt{o}}}\)\;
\end{algorithm}


%=========================================================
%=========================================================
%=========================================================
\subsection{Hierarchical abstraction and safety-ensured synthesis}
\label{ch06:sec:solution}

The baseline method in Section~\ref{ch06:sec:analysis} becomes expensive because
both temporal constraints are enforced directly on product automata built over
the original MDP \(\mathcal{M}\). To reduce this complexity, two feature-based
abstractions are introduced as semi-MDPs. One abstraction is constructed for
the safe-return specification, and the other is constructed for the nominal
task. Planning is then carried out on these abstractions rather than on the
full model.
The construction follows the semi-MDP and option framework of
\citet{sutton1999between}. Each abstract transition represents a motion policy
between two feature states. As in the baseline method, the return policy is
synthesized first, because its value function is required when building the
task-level abstraction. The structure is hierarchical: the return policy is
computed first and then used to construct the task-level abstraction for
outbound planning.


\subsubsection{Hierarchical synthesis of the return policy}
\label{ch06:subsec:safe-return-plan}

The return policy is synthesized first, together with the associated value
function that will later be used in the task abstraction. This requires a
feature-based abstraction tailored to the safe-return specification
\(\varphi_{\texttt{r}}\), followed by policy synthesis on the corresponding
abstract product.

(I) \emph{Return-level abstraction.}
Let
\(\mathcal{A}_{\textup{\texttt{r}}} \triangleq
(Q_{\textup{\texttt{r}}}, 2^{AP}, \delta_{\textup{\texttt{r}}},
q^{0}_{\textup{\texttt{r}}},
\mathrm{Acc}_{\mathcal{A}_{\textup{\texttt{r}}}})\)
be the deterministic Rabin automaton associated with
\(\varphi_{\texttt{r}}\). The effective labels for the return task are defined
by
\begin{equation}
\label{ch06:eq:effectve-safe}
\Theta_{\textup{\texttt{r}}} \triangleq
\left\{
\theta \in 2^{AP}
\;\middle|\;
\exists q \in Q_{\textup{\texttt{r}}}
\text{ such that }
\delta_{\textup{\texttt{r}}}(q,\theta) \neq q
\right\},
\end{equation}
which collects the labels that induce non-self transitions in
\(\mathcal{A}_{\textup{\texttt{r}}}\).
Using \(\Theta_{\textup{\texttt{r}}}\), a feature-based abstraction is
constructed as the labeled semi-MDP
\begin{equation}
\label{ch06:eq:abstract-mdp-safe}
\mathcal{M}'_{\textup{\texttt{r}}} \triangleq
\bigl(
X'_{\textup{\texttt{r}}},
U'_{\textup{\texttt{r}}},
p'_{\textup{\texttt{r}},D},
AP,
L'_{\textup{\texttt{r}}}
\bigr),
\end{equation}
where
$X'_{\textup{\texttt{r}}} \triangleq
\left\{
x \in X
\;\middle|\;
L(x) \in \Theta_{\textup{\texttt{r}}}
\right\}$,
and \(L'_{\textup{\texttt{r}}}\) is the restriction of the labeling map to
\(X'_{\textup{\texttt{r}}}\). Thus,
\(X'_{\textup{\texttt{r}}}\) contains only those states that are relevant to
the progress of \(\mathcal{A}_{\textup{\texttt{r}}}\). The abstract action set
\(U'_{\textup{\texttt{r}}}\) consists of macro actions, each of which
represents an underlying motion policy between two abstract states. The
transition probability
\(p'_{\textup{\texttt{r}},D} \colon
X'_{\textup{\texttt{r}}} \times
U'_{\textup{\texttt{r}}} \times
X'_{\textup{\texttt{r}}} \rightarrow [0,1]\)
is obtained as follows.

For each ordered pair
\((x_f,x_t) \in
X'_{\textup{\texttt{r}}} \times X'_{\textup{\texttt{r}}}\), the associated
macro action is denoted by
\((x_f,x_t) \in U'_{\textup{\texttt{r}}}\). A modified MDP
\(\mathcal{M}_{\textup{\texttt{r}}}(x_f,x_t)\) is constructed with source
state \(x_f\), while every state in
\(X'_{\textup{\texttt{r}}} \setminus \{x_f\}\) is treated as a sink state.
The option
\(\boldsymbol{\mu}'_{\textup{\texttt{r}}}(x_f,x_t)\) maximizes the probability
of reaching the target sink \(x_t\) \citep{sutton1999between}. Once this
option is fixed, the induced Markov chain defines the abstract transition
probability
\(p'_{\textup{\texttt{r}},D}
(x_f,\boldsymbol{\mu}'_{\textup{\texttt{r}}}(x_f,x_t),x_\ell)\)
for every \(x_\ell \in X'_{\textup{\texttt{r}}}\). Infeasible macro
transitions are omitted.
If the automaton
\(\mathcal{A}_{\textup{\texttt{r}}}\) contains \(N\) distinct relevant labels,
and at most \(M\) regions of interest are associated with each label in the
workspace, then \(\mathcal{M}'_{\textup{\texttt{r}}}\) has at most \(MN\)
states and fewer than \(M^2N^2\) abstract transitions. In large workspaces
with relatively few relevant regions, this abstraction can be substantially
smaller than the original MDP.


(II) \emph{Abstract return-policy synthesis.}
Given \(\mathcal{M}'_{\textup{\texttt{r}}}\) and the automaton
\(\mathcal{A}_{\textup{\texttt{r}}}\), their product is given by:
\begin{equation}
\label{ch06:eq:abstract-safe-product}
\mathcal{P}'_{\textup{\texttt{r}}} \triangleq
\mathcal{M}'_{\textup{\texttt{r}}}
\times
\mathcal{A}_{\textup{\texttt{r}}}
=
\bigl(
S'_{\textup{\texttt{r}}},
U'_{\textup{\texttt{r}}},
E'_{\textup{\texttt{r}}},
p'_{\textup{\texttt{r}},E},
s'^{0}_{\textup{\texttt{r}}},
\mathrm{Acc}'_{\textup{\texttt{r}}}
\bigr),
\end{equation}
where the notation is defined in the same way as in
Definition~\ref{ch06:def:product-o}. The transition probability
\(p'_{\textup{\texttt{r}},E}\) is induced from
\(p'_{\textup{\texttt{r}},D}\) in \eqref{ch06:eq:abstract-mdp-safe}.
The abstract return policy
\(\boldsymbol{\pi}'_{\textup{\texttt{r}}}\) and value function
\(v'^{\star}_{\textup{\texttt{r}}}\) are computed from
\(\mathcal{P}'_{\textup{\texttt{r}}}\). Since the safe-return specification is
sc-LTL, satisfaction reduces to reaching the union of AMECs, denoted by
\(S'_{\textup{\texttt{r}},\Xi}\).

\subsubsection{Hierarchical synthesis of the nominal task}
\label{ch06:subsec:task-plan}

After the return policy has been obtained, the nominal task is treated in the
same hierarchical spirit. The task abstraction differs from the return
abstraction in that it already incorporates the safe-return constraint through
the value function \(v'^{\star}_{\textup{\texttt{r}}}\).

(I) \emph{Safety-aware task abstraction.}
Let
\(\mathcal{A}_{\textup{\texttt{o}}}\) be the deterministic Rabin automaton
associated with \(\varphi\). The corresponding set of effective labels is
defined analogously to \eqref{ch06:eq:effectve-safe} and is denoted by
\(\Theta_{\textup{\texttt{o}}}\). The task-level abstraction is then given by
the semi-MDP below:
\begin{equation}
\label{ch06:eq:abstract-mdp-task}
\mathcal{M}'_{\textup{\texttt{o}}} \triangleq
\bigl(
X'_{\textup{\texttt{o}}},
U'_{\textup{\texttt{o}}},
p'_{\textup{\texttt{o}},D},
AP,
L'_{\textup{\texttt{o}}},
c'_{\textup{\texttt{o}},D}
\bigr),
\end{equation}
where the notation is analogous to
\eqref{ch06:eq:abstract-mdp-safe}. The main difference is that the abstract
transitions and abstract costs are computed under the additional safe-return
requirement.
For each ordered pair
\((x_f,x_t) \in
X'_{\textup{\texttt{o}}} \times X'_{\textup{\texttt{o}}}\), the associated
macro action is written as
\((x_f,x_t) \in U'_{\textup{\texttt{o}}}\). As before, a modified MDP
\(\mathcal{M}_{\textup{\texttt{o}}}(x_f,x_t)\) is constructed from the
original model. A motion policy
\(\boldsymbol{\mu}'_{\textup{\texttt{o}}}(x_f,x_t)\) is then synthesized to
maximize the probability of reaching \(x_t\), while also satisfying the
safe-return bound. By Theorem~\ref{ch06:theorem:safe-bound}, the latter can be
imposed through the constraint below:
\begin{equation}
\label{ch06:eq:abstract-task-lp}
\mathbb{E}_{\mathcal{M}_{\textup{\texttt{o}}}(x_f,x_t)}^{
\boldsymbol{\mu}'_{\textup{\texttt{o}}}}
\left[
\sum_{t=0}^{\infty}
v'^{\star}_{\textup{\texttt{r}}}(s^{\textup{\texttt{r}}}_t)
\right]
\geq
\chi_{\texttt{r}},
\end{equation}
where
\(s^{\textup{\texttt{r}}}_t \triangleq
\langle x_t,q^{0}_{\textup{\texttt{r}}}\rangle\)
is the product state in
\(\mathcal{P}'_{\textup{\texttt{r}}}\) associated with the MDP state
encountered at stage \(t\).
Once the option
\(\boldsymbol{\mu}'_{\textup{\texttt{o}}}(x_f,x_t)\) is fixed, the induced
Markov chain yields the abstract transition probabilities
\(p'_{\textup{\texttt{o}},D}\). The corresponding abstract cost is defined by:
\begin{equation}
\label{ch06:eq:task-edge-cost}
c'_{\textup{\texttt{o}},D}(x_f,x_t) \triangleq
\mathbb{E}_{\mathcal{M}_{\textup{\texttt{o}}}(x_f,x_t)}^{
\boldsymbol{\mu}'_{\textup{\texttt{o}}}}
\left[
\sum_{t=0}^{\infty}
c_D(x_t,u_t)
\right],
\end{equation}
that is, the expected cumulative cost of executing the option from \(x_f\)
until termination at one of the sink states. If the constrained linear program
is infeasible, then the transition \((x_f,x_t)\) is removed from
\(U'_{\textup{\texttt{o}}}\).
This local synthesis prioritizes reachability and safety over exact
macro-transition cost optimality. Hence
\(c'_{\textup{\texttt{o}},D}\) approximates the original objective in
\eqref{ch06:eq:objective}, but avoids global optimization on the full product.

(II) \emph{Abstract outbound-policy synthesis.}
Given the safety-ensured semi-MDP
\(\mathcal{M}'_{\textup{\texttt{o}}}\), the corresponding product automaton is
given by:
\begin{equation}
\label{ch06:eq:abstract-task-product}
\mathcal{P}'_{\textup{\texttt{o}}} \triangleq
\mathcal{M}'_{\textup{\texttt{o}}}
\times
\mathcal{A}_{\textup{\texttt{o}}},
\end{equation}
where the detailed notation is defined as in
Definition~\ref{ch06:def:product-o}, and the edge cost is induced by
\eqref{ch06:eq:task-edge-cost}.
Unlike the baseline product
\(\mathcal{P}_{\textup{\texttt{o}}}\), the abstract product
\(\mathcal{P}'_{\textup{\texttt{o}}}\) inherits the safe-return guarantee from
Lemma~\ref{ch06:lemma:safe}. Hence, no additional safe-return constraint is
imposed when synthesizing the outbound policy on
\(\mathcal{P}'_{\textup{\texttt{o}}}\). Since \(\varphi\) is a general LTL
formula, the policy consists of a prefix part
\(\boldsymbol{\pi}'_{\textup{\texttt{o}},\texttt{pre}}\), executed until an
accepting component is reached, and a suffix part
\(\boldsymbol{\pi}'_{\textup{\texttt{o}},\texttt{suf}}\), executed
thereafter. Following the framework of \citet{guo2018probabilistic}, the two
parts are obtained from the constrained optimization below:
\begin{subequations}
\label{ch06:eq:abstract-prod-task-lp}
\begin{align}
\min_{
\boldsymbol{\pi}'_{\textup{\texttt{o}},\texttt{pre}},
\boldsymbol{\pi}'_{\textup{\texttt{o}},\texttt{suf}}
}
\quad
&\textbf{Cost}_{\mathcal{P}'_{\textup{\texttt{o}}}}^{
\boldsymbol{\pi}'_{\textup{\texttt{o}},\texttt{suf}}},
\label{ch06:eq:abstract-prod-task-lp-objective} \\
\text{s.t.}\quad
&\textbf{Sat}_{\mathcal{P}'_{\textup{\texttt{o}}}}^{
\boldsymbol{\pi}'_{\textup{\texttt{o}},\texttt{pre}}}
\geq \chi_{\texttt{o}},
\label{ch06:eq:abstract-prod-task-lp-condition}
\end{align}
\end{subequations}
where the prefix policy ensures that the union of AMECs, denoted by
\(S'_{\textup{\texttt{o}},\Xi}\), is reached with probability at least
\(\chi_{\texttt{o}}\), while the suffix policy keeps the execution inside
\(S'_{\textup{\texttt{o}},\Xi}\) and minimizes the corresponding long-term
cost. The detailed linear program is omitted here.

\subsubsection{Algorithmic summary, online execution, and guarantee}
\label{ch06:subsec:summary}

Algorithm~\ref{ch06:alg:all} summarizes the procedure. It first constructs
the return abstraction, forms \(\mathcal{P}'_{\textup{\texttt{r}}}\), and
computes \(\boldsymbol{\pi}'_{\textup{\texttt{r}}}\) and
\(v'^{\star}_{\textup{\texttt{r}}}\). It then constructs the safety-aware
task abstraction, forms \(\mathcal{P}'_{\textup{\texttt{o}}}\), and
synthesizes the prefix and suffix outbound policies. Online execution selects
macro actions on the abstraction and executes the associated options on the
original MDP. Because options are stochastic, the high-level policy is
reapplied from the realized sink state. The outbound phase switches from prefix
to suffix control after entering \(S'_{\textup{\texttt{o}},\Xi}\); the
return phase is invoked only after a return request and terminates at
\(S'_{\textup{\texttt{r}},\Xi}\). Compared with the baseline, the
hierarchical method separates low-level option synthesis from high-level
temporal planning, which is the main source of computational savings.

\begin{algorithm}[t]
\caption{Hierarchical synthesis algorithm}
\label{ch06:alg:all}
\LinesNumbered
\KwIn{\(\mathcal{M}\), \((\varphi,\chi_{\texttt{o}})\), and
\((\varphi_{\texttt{r}},\chi_{\texttt{r}})\).}
\KwOut{\((\mathcal{M}'_{\textup{\texttt{r}}},
\boldsymbol{\mu}'_{\textup{\texttt{r}}}),
(\mathcal{P}'_{\textup{\texttt{r}}},
\boldsymbol{\pi}'_{\textup{\texttt{r}}},
v'^{\star}_{\textup{\texttt{r}}})\), and
\((\mathcal{M}'_{\textup{\texttt{o}}},
\boldsymbol{\mu}'_{\textup{\texttt{o}}}),
(\mathcal{P}'_{\textup{\texttt{o}}},
(\boldsymbol{\pi}'_{\textup{\texttt{o}},\texttt{pre}},
\boldsymbol{\pi}'_{\textup{\texttt{o}},\texttt{suf}}))\).}
Construct \(\mathcal{A}_{\textup{\texttt{r}}}\)\;
Build \(\mathcal{M}'_{\textup{\texttt{r}}}\) and its options
\(\boldsymbol{\mu}'_{\textup{\texttt{r}}}\)\;
Construct \(\mathcal{P}'_{\textup{\texttt{r}}}\) by
\eqref{ch06:eq:abstract-safe-product}\;
Synthesize \(\boldsymbol{\pi}'_{\textup{\texttt{r}}}\) and
\(v'^{\star}_{\textup{\texttt{r}}}\)\;
Construct \(\mathcal{A}_{\textup{\texttt{o}}}\)\;
Build \(\mathcal{M}'_{\textup{\texttt{o}}}\) and its options
\(\boldsymbol{\mu}'_{\textup{\texttt{o}}}\), using
\(v'^{\star}_{\textup{\texttt{r}}}\)\;
Construct \(\mathcal{P}'_{\textup{\texttt{o}}}\) by
\eqref{ch06:eq:abstract-task-product}\;
Synthesize
\((\boldsymbol{\pi}'_{\textup{\texttt{o}},\texttt{pre}},
\boldsymbol{\pi}'_{\textup{\texttt{o}},\texttt{suf}})\) from
\eqref{ch06:eq:abstract-prod-task-lp}\;
\end{algorithm}


\begin{lemma}
\label{ch06:lemma:safe}
The semi-MDP \(\mathcal{M}'_{\textup{\texttt{o}}}\) is safety-ensured.
More precisely, for every transition
\((x_f,x_t) \in U'_{\textup{\texttt{o}}}\), the associated motion policy
\(\boldsymbol{\mu}'_{\textup{\texttt{o}}}(x_f,x_t)\), together with the
return policy \(\boldsymbol{\mu}'_{\textup{\texttt{r}}}\), satisfies the
safe-return constraint in \eqref{ch06:eq:safe}.
\end{lemma}

\begin{proof}
For each macro transition \((x_f,x_t)\), the motion policy
\(\boldsymbol{\mu}'_{\textup{\texttt{o}}}(x_f,x_t)\) is synthesized subject
to \eqref{ch06:eq:abstract-task-lp}. By
Theorem~\ref{ch06:theorem:safe-bound}, the left-hand side of
\eqref{ch06:eq:abstract-task-lp} is exactly the safe-return probability
associated with the option execution and the return policy. Hence, every
admissible transition in \(\mathcal{M}'_{\textup{\texttt{o}}}\) preserves the
probability bound \(\chi_{\texttt{r}}\). The entire abstraction is therefore
safety-ensured.
\end{proof}


\begin{theorem}
\label{ch06:thm:safety-guarantee}
The hierarchical policies
\(\boldsymbol{\pi}'_{\textup{\texttt{r}}}\) and
\(\boldsymbol{\pi}'_{\textup{\texttt{o}}}\) satisfy the task and safe-return
constraints in \eqref{ch06:eq:objective}. Moreover, the optimization in
\eqref{ch06:eq:abstract-prod-task-lp} minimizes the cost induced on the
abstract product \(\mathcal{P}'_{\textup{\texttt{o}}}\).
\end{theorem}

\begin{proof}
By Lemma~\ref{ch06:lemma:safe}, every macro transition in
\(\mathcal{M}'_{\textup{\texttt{o}}}\) satisfies the safe-return bound
\(\chi_{\texttt{r}}\) together with the return policy
\(\boldsymbol{\mu}'_{\textup{\texttt{r}}}\). Hence, any execution generated by
the outbound policy on \(\mathcal{P}'_{\textup{\texttt{o}}}\) preserves the
safe-return constraint in \eqref{ch06:eq:objective}.
For the task requirement, the prefix-suffix synthesis on
\(\mathcal{P}'_{\textup{\texttt{o}}}\) follows the same automata-based
construction used by \citet{guo2018probabilistic}. In particular, the prefix
policy ensures that the union of AMECs \(S'_{\textup{\texttt{o}},\Xi}\) is
reached with probability at least \(\chi_{\texttt{o}}\), and the suffix policy
maintains the execution within the accepting component. Therefore, the
task-satisfaction bound in \eqref{ch06:eq:objective} is satisfied.
Finally, \eqref{ch06:eq:abstract-prod-task-lp-objective} minimizes the cost
defined on the abstract product \(\mathcal{P}'_{\textup{\texttt{o}}}\).
Because the local options used to construct
\(\mathcal{M}'_{\textup{\texttt{o}}}\) prioritize reachability and safety
rather than exact transition-cost minimization, this hierarchical cost is an
approximation of the original cost in \eqref{ch06:eq:objective}. The
constraint satisfaction, however, is preserved.
\end{proof}
